\scanpage{55}
\uline{Catégories $\operatorname{Ind}$}$_\pi(C)$ ($\deleted{\unsure{\pi}} C$ petite) et $\deleted{\operatorname{Ind}_\pi(C)_{\pi'}\unsure{(\mathscr E)}}\operatorname{Fil}^{\pi}(\operatorname{Ind}_{\pi}(C))$

\uline{Th.} $E$ \inserted{$\mathcal U$-}catégorie [stable par petites $\varinjlim$ ? \inserted{ou satisfaisant seulement $L_\pi$ ?}]. Conditions équivalentes

a) $E$ $\pi$-accessible

b) $E$ équivalente à une catégorie $\operatorname{Ind}_\pi(C)$, $C$ petite

c) La sous-catégorie $E_\pi$ de $E$ formée des objets $\pi$-accessibles est équiv. à une petite cat. et elle est strictement génératrice.

\deleted{et de même une filtration ordinale, si $\pi'\geq\pi$ \unsure{cf hypothèse} $\pi'\geq\pi$.}

Sous ces conditions, \deleted{$E$ admet une filtration \unsure{ordinale|cardinale}, $(\unsure{E^\circ_\pi}(E_\pi))_{\pi'\geq\pi}$} on a $E_{\pi'}\simeq\unsure{\operatorname{Ind}}(\operatorname{Ind}_{\pi}(E_\pi)_{\pi'})$, et les $(E_{\pi'})_{\pi'\geq\pi}$ forment une filtration ordinale
\deleted{\unsure{de filtration}}.

Préciser la filtration ordinale \deleted{\unsure{de|des}} $(\underline{\operatorname{Fil}}^{\pi})_{\pi\geq\pi_0}$ \unsure{---}
\[
 \operatorname{Fil}^{\pi'}(E)=E_{\pi'}\inserted{{}=\operatorname{Ind}_{\pi}(C)_{\pi'}}\quad\text{pour $\pi'$ grand,}
\]

\uline{Cor} Soient $E$ et $F$ deux catégories accessibles, \inserted{donc $E\simeq\operatorname{Ind}_\pi(E_\pi)$, $F\simeq\operatorname{Ind}_\pi(F_\pi)$ pour $\pi\geq\pi_0$ \deleted{gd}} \deleted{Alors} $f:E\to F$ un foncteur. \deleted{Alors} Conditions équiv..

a) $f$ est \deleted{$\pi$-}accessible \deleted{\unsure{c)}} b) $\exists\ \pi$ tel que $f\simeq\operatorname{Ind}_\pi(f_\pi)$, où $f_\pi:E_\pi\to F_\pi$ un foncteur. \deleted{[et pour un tel $\pi$, $f$ est $\pi$-accessible]. Preuve que} [Pr. v. que pr $\pi$ canonique, il $f$ est \unsure{''} que $f$ soit $\pi$-accessible et envoie $E_\pi$ dans $F_\pi$]. c) $\exists$ une petite sous-catégorie pleine $C$ de $E$, $i:C\to E$ telle que $f$ soit dans l'image essentielle de $i_!^F$ :
\deleted{b) Pour un tel $\pi$, $\unsure{i|f_i}$ est} $\underline{\operatorname{Hom}}(C,F)\to\underline{\operatorname{Hom}}(E,F)$.

\uline{Corollaire} Supposons $f$ $\pi$-accessible. Alors \deleted{\unsure{il y a}} $\exists\ \pi_0$ tel que
\[
 \pi\geq\pi_0\Longrightarrow f(E_\pi)\subset F_\pi .
\]
\margin{\uline{Corollaire} \inserted{$E$ $\pi$-accessible, $F$ qlc. Alors} $\underline{\operatorname{Hom}}(E,F)_\pi\overset{\deleted{\unsure{\sim}}}{\longrightarrow}\underline{\operatorname{Hom}}(E_\pi,F)$ pl. fid, c'est une équivalence de catégories \deleted{\unsure{ssi}} si $F$ satisfait $L_\pi$}

\uline{Th} Soit $C$ une catégorie accessible \inserted{avec petites $\varinjlim$.} \deleted{\unsure{$f$:}} Alors

a) (Pour mémoire) les foncteurs contravariants $C^\circ\to(\mathrm{Ens})$ repr. sont ceux qui commutent aux \inserted{petites} $\varinjlim$\deleted{\unsure{---}}

\margin{8.13.2 $(+8.12.8.1)$}

b) Un foncteur covariant \deleted{\unsure{en}} $C\to(\mathrm{Ens})$ est proreprésentable \inserted{(resp.)} \deleted{ssi $\unsure{\Pi|\pi}$ est accessible et} \inserted{ssi il est} exact à gauche (resp. commute aux petites $\varprojlim$) et \uline{accessible}

\uline{Corollaire} Soit $f:C\to C'$ un foncteur, avec $C$ accessible \inserted{stable par petites $\varinjlim$}.

a) $f$ admet un adjoint à droite ssi $f$ commute aux \inserted{petites} $\varinjlim$

b) $f$ admet un \inserted{pro}adjoint \deleted{\unsure{en}} (resp. un adjoint) à gauche ssi il est exact à \unsure{g} (resp. commute aux \inserted{petites} $\varprojlim$) \deleted{que} et est \uline{accessible}.

\uline{Th.} a) Soit $E\to B$ une catégorie fibrée, avec $B$ équiv. à une \deleted{\unsure{petit|petite}} \inserted{petite} catégorie, les catégories fibres accessibles, les foncteurs changement de base accessibles. Alors $\varprojlim E/B$ et $\underline{\operatorname{Hom}}_B(B,E)$ sont accessibles. En particulier, si $C$ accessible et $I$ équiv. à une petite catégorie, $\underline{\operatorname{Hom}}(I,C)$ est accessible, ainsi que la sous-catégorie pleine des foncteurs qui transforment tte flèche en un iso.

\margin{9.22}

b) Même hyp. sur $E,B$. Alors $\varinjlim E/B^\circ$ est accessible.
\margin{à \unsure{voir|revoir|établir|vérifier} $\longrightarrow$}% slanted
% archive sheet number: 54
